% GENERATED FILE. Edit canonical lessons, then run npm run generate:books.
% canonical-source-hash: fnv1a64:d0acb66d2f5be5b9
% canonical-lessons: FR-C500-belote, FR-C500-distribuer, FR-C500-parier, FR-C500-pari, FR-C500-hasard

\chapter{Cartes et paris}
\label{ch:fr-cartes-et-paris}
\hlchaptermodality{500}

\begin{chapteropening}
\textbf{By the end of this chapter:} \emph{I can talk about card games and betting.}

Complete the last lesson of chapter 500: I can talk about card games and betting.
\end{chapteropening}

\section[la belote]{la belote — belote (a card game)}

\label{lesson:FR-C500-belote}



\begin{quote}
Before the new one: say the French for chessboard, then the French for to cheat.
\end{quote}

\subsection*{You'll want to know: la belote}
\textbf{la belote} — \textquotedblleft{}belote (a card game)\textquotedblright{}.

\begin{quote}
Le dimanche, mes oncles jouent à la belote.
\end{quote}

\begin{quote}
\emph{On Sundays my uncles play belote.}
\end{quote}

A very popular card game in France, played in pairs with 32 cards.

\begin{cousinweb}
The origin of \emph{la belote} is not settled; the game took its name in France in the early twentieth century.
\end{cousinweb}

\begin{grammarlens}[title={What you've built}]
One more word for telling, explaining and describing at B1.
\end{grammarlens}

\subsection*{Your turn}
\begin{itemize}
\raggedright
  \item \emph{Say it:} \emph{la belote}
  \item \emph{Say it:} \emph{la belote}, once more
  \item \emph{Say it:} a sentence of your own with \emph{la belote}
  \item \emph{Recall:} say the French for chessboard, then the French for to cheat, then say \emph{la belote} again
\end{itemize}

\begin{tcolorbox}[breakable,colback=teal!4,colframe=teal!35!black,title={Before you move on}]
What does la belote mean? (\textquotedblleft{}Belote (a card game)\textquotedblright{}.) Say it once more.
\end{tcolorbox}
\section[distribuer]{distribuer — to deal (cards); to hand out}

\label{lesson:FR-C500-distribuer}



\begin{quote}
Before the new one: say the French for to cheat, then the French for belote.
\end{quote}

\subsection*{You'll want to know: distribuer}
\textbf{distribuer} — \textquotedblleft{}to deal (cards); to hand out\textquotedblright{}.

\begin{quote}
C'est à ton tour de distribuer les cartes.
\end{quote}

\begin{quote}
\emph{It is your turn to deal the cards.}
\end{quote}

Also to hand out: \emph{distribuer des tracts}.

\begin{cousinweb}
\emph{distribuer} is a learned borrowing from Latin \emph{distribuere} \textquotedblleft{}to divide out\textquotedblright{}, built on \emph{tribuere} \textquotedblleft{}to assign\textquotedblright{} (Lewis \& Short).
\end{cousinweb}

\begin{grammarlens}[title={What you've built}]
One more word for telling, explaining and describing at B1.
\end{grammarlens}

\subsection*{Your turn}
\begin{itemize}
\raggedright
  \item \emph{Say it:} \emph{distribuer}
  \item \emph{Say it:} \emph{distribuer}, once more
  \item \emph{Say it:} a sentence of your own with \emph{distribuer}
  \item \emph{Recall:} say the French for to cheat, then the French for belote, then say \emph{distribuer} again
\end{itemize}

\begin{tcolorbox}[breakable,colback=teal!4,colframe=teal!35!black,title={Before you move on}]
What does distribuer mean? (\textquotedblleft{}To deal (cards); to hand out\textquotedblright{}.) Say it once more.
\end{tcolorbox}
\section[parier]{parier — to bet}

\label{lesson:FR-C500-parier}



\begin{quote}
Before the new one: say the French for belote, then the French for to deal.
\end{quote}

\subsection*{You'll want to know: parier}
\textbf{parier} — \textquotedblleft{}to bet\textquotedblright{}.

\begin{quote}
Il a parié dix euros sur son équipe.
\end{quote}

\begin{quote}
\emph{He bet ten euros on his team.}
\end{quote}

\emph{Je parie que…} — I bet that…

\begin{cousinweb}
\emph{parier} comes from Late Latin \emph{pariāre} \textquotedblleft{}to make equal, to match\textquotedblright{}, built on \emph{pār} \textquotedblleft{}equal\textquotedblright{} (Lewis \& Short): two sides put up equal stakes.
\end{cousinweb}

\begin{grammarlens}[title={What you've built}]
One more word for telling, explaining and describing at B1.
\end{grammarlens}

\subsection*{Your turn}
\begin{itemize}
\raggedright
  \item \emph{Say it:} \emph{parier}
  \item \emph{Say it:} \emph{parier}, once more
  \item \emph{Say it:} a sentence of your own with \emph{parier}
  \item \emph{Recall:} say the French for belote, then the French for to deal, then say \emph{parier} again
\end{itemize}

\begin{tcolorbox}[breakable,colback=teal!4,colframe=teal!35!black,title={Before you move on}]
What does parier mean? (\textquotedblleft{}To bet\textquotedblright{}.) Say it once more.
\end{tcolorbox}
\section[le pari]{le pari — bet}

\label{lesson:FR-C500-pari}



\begin{quote}
Before the new one: say the French for to deal, then the French for to bet.
\end{quote}

\subsection*{You'll want to know: le pari}
\textbf{le pari} — \textquotedblleft{}bet\textquotedblright{}.

\begin{quote}
J'ai perdu mon pari avec mon frère.
\end{quote}

\begin{quote}
\emph{I lost my bet with my brother.}
\end{quote}

\emph{Faire un pari} — to make a bet.

\begin{cousinweb}
\emph{le pari} is built on \emph{parier} \textquotedblleft{}to bet\textquotedblright{} (EtymDB), which comes from Late Latin \emph{pariāre} \textquotedblleft{}to make equal\textquotedblright{}, from \emph{pār} (Lewis \& Short).
\end{cousinweb}

\begin{grammarlens}[title={What you've built}]
One more word for telling, explaining and describing at B1.
\end{grammarlens}

\subsection*{Your turn}
\begin{itemize}
\raggedright
  \item \emph{Say it:} \emph{le pari}
  \item \emph{Say it:} \emph{le pari}, once more
  \item \emph{Say it:} a sentence of your own with \emph{le pari}
  \item \emph{Recall:} say the French for to deal, then the French for to bet, then say \emph{le pari} again
\end{itemize}

\begin{tcolorbox}[breakable,colback=teal!4,colframe=teal!35!black,title={Before you move on}]
What does le pari mean? (\textquotedblleft{}Bet\textquotedblright{}.) Say it once more.
\end{tcolorbox}
\section[le hasard]{le hasard — chance, luck}

\label{lesson:FR-C500-hasard}



\begin{quote}
Before the new one: say the French for to bet, then the French for bet.
\end{quote}

\subsection*{You'll want to know: le hasard}
\textbf{le hasard} — \textquotedblleft{}chance, luck\textquotedblright{}.

\begin{quote}
Les jeux de hasard sont interdits aux mineurs.
\end{quote}

\begin{quote}
\emph{Games of chance are forbidden to minors.}
\end{quote}

\emph{Par hasard} — by chance. The \emph{h} is aspirated: \emph{le hasard}.

\begin{cousinweb}
\emph{le hasard} comes from Arabic \emph{az-zahr} \textquotedblleft{}the die\textquotedblright{}, through Spanish \emph{azar} (EtymDB): first a game of dice, then chance itself.
\end{cousinweb}

\begin{grammarlens}[title={What you've built}]
One more word for telling, explaining and describing at B1.
\end{grammarlens}

\subsection*{Your turn}
\begin{itemize}
\raggedright
  \item \emph{Say it:} \emph{le hasard}
  \item \emph{Say it:} \emph{le hasard}, once more
  \item \emph{Say it:} a sentence of your own with \emph{le hasard}
  \item \emph{Recall:} say the French for to bet, then the French for bet, then say \emph{le hasard} again
\end{itemize}

\begin{tcolorbox}[breakable,colback=teal!4,colframe=teal!35!black,title={Before you move on}]
What does le hasard mean? (\textquotedblleft{}Chance, luck\textquotedblright{}.) Say it once more.
\end{tcolorbox}
